% Chapter 9 — keep it running: orchestration, quality, operations.

\gmenter{operate}
\note{Chapter 9: the machine around all of it --- the dashed box on the map. Three stops: the
monthly chain, the machine it runs on and the rules that made it hold, and how it is checked.}

\gmgo{operate-chain}{The monthly chain}
\begin{frame}{Every month, the chain redoes only what changed}
  \centering
  \begin{tikzpicture}[x=1cm, y=1cm,
      st/.style={circle, minimum size=1.15cm, inner sep=0pt, line width=1.6pt, fill=gmSurface,
                 font=\fontsize{15}{15}\selectfont},
      lb/.style={font=\small, text=gmInk, align=center, text width=2.05cm},
      go/.style={-{Stealth[length=2.6mm]}, line width=1.2pt, draw=gmBase, shorten <=1mm, shorten >=1mm}]
    \foreach \i/\c/\ic/\t in {0/collect/\faSitemap/find pages, 1/collect/\faGlobe/fetch pages,
                              2/read/\faFile*/read pages, 3/link/\faLink/link the names,
                              4/link/\faGavel/keep what matters, 5/events/\faStream/tell the story} {
      \node[st, draw=gm@col@\c, text=gm@col@\c] (a\i) at (\i*2.4,1.75) {\ic};
      \node[lb, below=1mm of a\i] {\t};}
    \foreach \i/\c/\ic/\t in {0/graph/\faProjectDiagram/build \& publish, 1/link/\faBook/Wikidata facts,
                              2/picture/\faImage/read the meme's own image, 3/picture/\faEye/read the templates,
                              4/picture/\faImages/find templates} {
      \node[st, draw=gm@col@\c, text=gm@col@\c] (b\i) at (\i*2.4+1.2,-1.45) {\ic};
      \node[lb, below=1mm of b\i] {\t};}
    \foreach \i [count=\j] in {0,...,4} \draw[go] (a\i) -- (a\j);
    \draw[go] (a5.east) to[out=0, in=0, looseness=1.4] (b4.east);
    \foreach \i [count=\j] in {0,...,3} \draw[go] (b\j) -- (b\i);
    \coordinate (lk) at ($(a3.south)+(0,-1.2)$);
    \draw[gm@col@link, line width=1.2pt, dash pattern=on 3pt off 2.5pt, -{Stealth[length=2.4mm]}] (lk) to[out=-90, in=90] (b3.north);
    \draw[gm@col@link, line width=1.2pt, dash pattern=on 3pt off 2.5pt, -{Stealth[length=2.4mm]}] (lk) to[out=-90, in=90] (b2.north);
    \node[fill=gmSurface, font=\small\bfseries, text=gm@col@link!85!gmInk, inner sep=1.5pt, anchor=west] at ($(lk)+(.15,-.05)$) {the same linker};
    \node[font=\normalsize\bfseries, text=gm@col@operate, anchor=south] at ($(a0.north)+(0,.1)$) {1st of the month};
  \end{tikzpicture}
  \vskip4mm
  {\large Each step does \textbf{only what is new or changed} --- and can stop and restart safely.}
  \note{
  \begin{itemize}
    \item Layman version: every month a robot repeats the whole journey we just made with Doge ---
      but only for the pages that are new or changed. Doge, unchanged, is not read again.
    \item The chain, as the stages trigger each other: find pages, fetch them, read them, link the
      names in the text, keep what matters, tell the story, find templates, read them, read each
      meme's own image, fetch the Wikidata facts of everything linked, build and publish the
      graph. Twelve Airflow DAGs; each starts the next.
    \item Why linking is this early: linking the text needs only the parsed page. The events are not
      linked (their people and places stay the page's words), and the pictures come later --- their
      names go through the same linker, which prefers the items the page's own text already
      linked. So the text has to be linked first.
    \item If anything fails half-way, running it again finishes the job without redoing or losing
      work. That is the property everything is built around.
  \end{itemize}}
\end{frame}

\gmgo{operate-machine}{The machine}
\begin{frame}{The machine, and the rules that made it hold}
  \begin{columns}[c]
    \column{.58\textwidth}\centering
      \resizebox{\linewidth}{!}{\begin{tikzpicture}[x=1cm, y=1cm,
          bx/.style={rounded corners=1.6mm, draw=gmBase, line width=1pt, fill=gmSurface, align=center, font=\normalsize,
                     minimum height=1.1cm, inner sep=2pt},
          go/.style={-{Stealth[length=2mm]}, line width=.9pt, draw=gmBase},
          grp/.style={rounded corners=2.5mm, draw=gmHair, line width=.8pt, fill=gmTint}]
        \draw[grp] (-0.4,-0.7) rectangle (5.2,3.4);  \node[font=\normalsize\bfseries, text=gmInk2, anchor=north west] at (-0.3,3.35) {Airflow 3: \FDags\ DAGs};
        \draw[grp] (5.6,-0.7) rectangle (10.3,3.4); \node[font=\normalsize\bfseries, text=gmInk2, anchor=north west] at (5.7,3.35) {stores};
        \node[bx, draw=gm@col@operate, minimum width=2.2cm] (sch) at (1.1,2.2) {scheduler};
        \node[bx, draw=gm@col@operate, minimum width=2.2cm] (wk) at (3.7,2.2) {workers};
        \node[bx, minimum width=2.2cm] (pg) at (1.1,0.4) {run history};
        \node[bx, minimum width=2.2cm] (dag) at (3.7,0.4) {stage code};
        \node[bx, draw=gm@col@read, minimum width=2.0cm] (mo) at (6.9,2.2) {MongoDB};
        \node[bx, draw=gm@col@graph, minimum width=2.0cm] (n4) at (9.1,2.2) {Neo4j};
        \node[bx, draw=gm@col@graph, minimum width=2.0cm] (fu) at (9.1,0.4) {Fuseki};
        \node[bx, minimum width=2.0cm] (fi) at (6.9,0.4) {files};
        \node[bx, draw=gm@col@collect, minimum width=3.2cm] (web) at (2.4,4.75) {Know Your Meme, imgflip};
        \node[bx, draw=gm@col@events, minimum width=3.2cm] (gpu) at (7.9,4.75) {the lab's GPUs};
        \node[bx, draw=gm@col@numbers, minimum width=5.2cm] (px) at (4.9,-2.0) {one door: port 8080};
        \draw[go] (web) -- (wk);  \draw[go] (wk) -- (gpu);
        \draw[go] (wk) -- (mo);   \draw[go] (wk.south east) -- (fi.north west);
        \draw[go] (mo) -- (n4);   \draw[go] (fi) -- (fu);
        \draw[go] (px.north) -- ++(0,.8);
      \end{tikzpicture}}
    \column{.4\textwidth}
      \begin{tikzpicture}[x=1cm, y=-1.3cm,
          ic/.style={circle, minimum size=.9cm, fill=gm@col@operate, text=white, font=\fontsize{12}{12}\selectfont, inner sep=0pt},
          tx/.style={anchor=west, text width=4.5cm, font=\normalsize, text=gmInk, align=left}]
        \node[ic] at (0,0) {\faStamp};     \node[tx] at (.6,0) {\textbf{Stamps}\\[-.4mm]{\small\color{gmInk2}redo only what changed}};
        \node[ic] at (0,1) {\faInbox};     \node[tx] at (.6,1) {\textbf{Dead letters}\\[-.4mm]{\small\color{gmInk2}failures kept, with the reason}};
        \node[ic] at (0,2) {\faMicrochip}; \node[tx] at (.6,2) {\textbf{Models on record}\\[-.4mm]{\small\color{gmInk2}every output names its model}};
        \node[ic] at (0,3) {\faChartLine}; \node[tx] at (.6,3) {\textbf{A summary per run}\\[-.4mm]{\small\color{gmInk2}read by the dashboard}};
      \end{tikzpicture}
  \end{columns}
  \note{
  \begin{itemize}
    \item 15 services in Docker Compose on one server. Airflow 3 orchestrates 12 DAGs: a scheduler,
      Celery workers with Redis, Postgres for Airflow's own history. The stage code is layered: a
      pure library, a store module that owns its collection, the DAG; a DAG never issues a query,
      a library never touches a database.
    \item MongoDB holds every stage's output; the graph goes to Neo4j and Fuseki, plus files.
      Models run on the lab's GPU servers (Ollama behind Open WebUI): ministral-3:14b for text,
      qwen3-vl:32b for images. One door: the campus firewall lets only port 8080 through, so one
      proxy serves the dashboard, Neo4j Browser, the SPARQL endpoint and the admin tools.
    \item Stamps, not timestamps: each output records the versions that produced it --- parser,
      prompt, schema, model, lexicon, linker. Changing the event prompt re-extracts events; it
      does not re-scrape pages.
    \item Dead letters: parse, event and template failures are stored with their reason, and not
      retried until something relevant changes.
    \item Models on record: every output names the model that produced it, so the graph never
      mixes readings without saying so. The shared GPUs shaped everything: reading 36,673 stories
      and 26,868 pictures took days, one call at a time, so every step had to survive
      interruption and resume where it stopped.
  \end{itemize}}
\end{frame}

\gmgo{operate-checked}{Checked and versioned}
\begin{frame}{Checked, measured, versioned}
  \begin{columns}[t]
    \column{.3\textwidth}
      \gmbig{\FTests}{automated tests}\par\vskip3mm
      \gmbig{\FGaps}{technical gaps, written down}\par\vskip3mm
      \gmbig{\FCommits}{commits since \FStart}
    \column{.3\textwidth}
      {\normalsize\color{gmInk2}every build, measured}\par\vskip3mm
      {\large\begin{tabular}{@{}rl@{}}
        \textbf{\NIsolated}       & isolated nodes\\[1mm]
        \textbf{\NDuplicateEdges} & duplicate edges\\[1mm]
        \textbf{0}                & cycles in a series\\[1mm]
        \textbf{=}                & RDF, derived twice\\
      \end{tabular}}
    \column{.36\textwidth}
      {\normalsize\color{gmInk2}a graph with a version}\par\vskip3mm
      {\normalsize\begin{tabular}{@{}rl@{}}
        \textbf{4.0} & nodes: only what is shared\\
        \textbf{5.0} & origins, folded tags\\
        \textbf{6.0} & events\\
        \textbf{6.1} & Wikidata links\\
        \textbf{6.4} & templates\\
        \textbf{6.5} & curated links\\
        \textbf{7.0} & siblings; a rename\\
        \textbf{7.1} & images; Wikidata's facts\\
      \end{tabular}}
  \end{columns}
  \note{
  \begin{itemize}
    \item 1,175 tests, about 9k of the 40k lines of Python. Parsers are tested on saved pages;
      model-facing code on recorded answers; the chain itself is tested (who triggers whom, with
      which parameters). Every model layer is reviewed in a loop --- sample, read, fix, read again
      --- on contact sheets, blind audits and held-out sets never used for tuning.
    \item 14 technical gaps logged in the repository: each says what is weak, what it costs, and the
      fix. Several chapters today closed one.
    \item After every publish, kg/metrics.py measures the IMKG-comparable core: no isolated node, no
      duplicate edge, no cycle in any series chain; and the RDF derived independently through the
      mapping equals the published RDF. Worth knowing: 26\% of series parents are pages not in
      the corpus (they exist as stubs); 18.8\% of entries have no entry type on KYM.
    \item Versions follow semantic versioning: a renamed term breaks queries, so 7.0.0 is major
      (relatesToMeme became citesMediaFrame); a new layer only adds, so 7.1.0 is minor. The build
      compares every input's stamp with the published build's: nothing changed, nothing rebuilt.
    \item A version worth telling: 4.0.0 decided what deserves to be a node --- only what things
      share. The graph went from 985,481 to 476,794 nodes with no data lost.
    \item 96 commits since 24 June: the whole project in about three and a half months.
  \end{itemize}}
\end{frame}
